\begin{studiumworked}{M26. Dos etapas: separación sin impulso · Intermedio}
\textbf{Enunciado.} Un vehículo ideal lleva carga de 65 kg. La primera etapa tiene masa seca $126\,\mathrm{kg}$ y propulsante $650\,\mathrm{kg}$, con $c_1=2450$ m/s. La segunda tiene 40 kg secos y $186\,\mathrm{kg}$ de propulsante, con $c_2=2900$ m/s. Calcule masas de inicio y fin de cada combustión y delta-v total. Cuantifique el error de retirar la masa seca de primera etapa antes de terminar su combustión.

\textbf{Datos.} Seca primera: $126\,\mathrm{kg}$; Propulsante primera: $650\,\mathrm{kg}$; Seca segunda: $40\,\mathrm{kg}$; Propulsante segunda: $186\,\mathrm{kg}$; Carga: $65\,\mathrm{kg}$.

\textbf{Incógnitas.} Masas anidadas y delta-v por etapa.

\textbf{Hipótesis.} Etapas en serie, sin pérdidas ni reservas. Separación sin velocidad relativa; velocidades efectivas constantes.

\textbf{Modelo físico.} Tsiolkovski por intervalos y contabilidad de masas.
\[m_{01}=s_1+p_1+s_2+p_2+m_L,\quad m_{f1}=s_1+s_2+p_2+m_L\]
\[m_{02}=s_2+p_2+m_L,\quad m_{f2}=s_2+m_L\]
\[\Delta v=c_1\ln(m_{01}/m_{f1})+c_2\ln(m_{02}/m_{f2})\]
\textbf{Desarrollo.}
\begin{enumerate}
\item Antes de la primera combustión hay $1067\,\mathrm{kg}$. Al agotarla quedan $417\,\mathrm{kg}$, incluidos $126\,\mathrm{kg}$ secos de esa etapa. Separarla deja $291\,\mathrm{kg}$, que es la masa inicial de segunda.
\item Al terminar la segunda quedan $105\,\mathrm{kg}$. Los logaritmos son $2450\ln(1067/417)=2301.824$ m/s y $2900\ln(291/105)=2956.152$ m/s.
\item Sumar da $5257.977\,\mathrm{m/s}$. Usar $291\,\mathrm{kg}$ en el denominador del primer logaritmo daría $6139.396\,\mathrm{m/s}$, un exceso de $881.4192\,\mathrm{m/s}$. Ese exceso atribuye a la separación un efecto propulsivo que el modelo no contiene.
\end{enumerate}
\textbf{Resultados.} $\Delta v_1=2301.824$ m/s; $\Delta v_2=2956.152$ m/s; total $5257.977$ m/s.

\textbf{Interpretación.} La segunda etapa acelera una masa inferior gracias al desprendimiento previo. La velocidad no salta al separar masa seca sin impulso.

\textbf{Comprobación.} Las diferencias de masa de cada combustión son exactamente sus propulsantes. La diferencia entre fin de primera y comienzo de segunda es exactamente $126\,\mathrm{kg}$.

\textbf{Límites.} No representa consumo paralelo, trayectoria, masa de interetapas no incluida ni variación de prestaciones con ambiente.
\end{studiumworked}
